% !TeX root = ../动手学强化学习笔记.tex
% 由主文档通过 \input 载入；请编译主文档。
\setcounter{chapter}{2} % 保留原书章号。
\chapter{部分习题解答与进一步学习}
\label{app:rl-exercises}

\section{基础计算题解答}

第\ref{ch:rl-introduction}章的路线阈值由 $-1+5\gamma=2$ 得 $\gamma=3/5$；低于阈值选立即奖励，高于阈值选延迟路线，等号时并列。若第二步以概率 $0.8$ 得 $5$，期望奖励为 $4$，阈值变为 $3/4$。两者比较的都是期望折扣回报。

第\ref{ch:rl-math}章奖励序列 $1,5,0,2$ 的均值依次为 $1,3,2,2$；常数步长 $1/2$、初值零时依次为 $0.5,2.75,1.375,1.6875$。霍夫丁界要求 $2e^{-0.02n}\leq0.05$，即 $n\geq\log40/0.02$，取整数 $185$ 足够；这是充分样本量，不宣称必要。第\ref{ch:rl-bandits}章三动作的选择概率为 $(0.1,0.45,0.45)$，伪遗憾题的答案为 $0.2\times200=40$；两个相同连续后验因交换对称性各有 $1/2$ 的选择概率。

第\ref{ch:rl-mdp}章可取 $\Pr(T=n)=1/[n(n+1)]$，$n\geq1$。总概率通过裂项求和等于一，而 $\mathbb E[T]=\sum_{n\geq1}1/(n+1)=\infty$。第\ref{ch:rl-bellman}章两状态价值为 $v_2=2/(1-\gamma)$、$v_1=1/[(1-\gamma)(1-\gamma/2)]$；奖励为正且持续时，折扣趋近一使价值发散，不代表平均奖励也发散。

第\ref{ch:rl-dp}章的一状态价值误差为 $20(0.9)^k$，小于 $0.1$ 需要 $k>\log(0.005)/\log(0.9)$，故 $k=51$ 次足够。贪心策略从对零表第一次比较动作时就已选对。第\ref{ch:rl-mc}章短回合的回报按逆序为 $2,2,1$，故状态 $s$ 的首次访问估计为 $1$，每次访问估计为 $1.5$。

\section{算法计算题解答}

第\ref{ch:rl-td}章若 Sarsa 的下一动作改为价值 $4$ 的动作，其目标变为 $4.6$，新估计为 $3.16$；Q-learning 与期望 Sarsa 不依赖这次实际抽到哪个后继动作，结果不变。第\ref{ch:rl-multistep}章访问 $s,u,s$ 时，累计迹依次为 $(1,0),(1/2,1),(5/4,1/2)$，替换迹最后为 $(1,1/2)$，其中坐标按 $(s,u)$ 排列。

第\ref{ch:rl-offpolicy}章目标动作期望为 $0.4\times1+0.6\times3=2.2$；行为加权计算为 $0.8(0.4/0.8)1+0.2(0.6/0.2)3=2.2$。第\ref{ch:rl-approximation}章常数近似在均匀权重下为 $2$，在权重 $(0.9,0.1)$ 下为 $1.2$，均由加权平方损失求导得到。

第\ref{ch:rl-actor-critic}章末尾改为截断、后继价值 $0.6$ 后，最后残差为 $1+0.9\times0.6-0.2=1.34$，按时间正序的优势为 $(0.03235,0.383,1.34)$。截断处 $b=1,c=0$。第\ref{ch:rl-continuous}章软最优动作概率为 $(1/4,3/4)$，软价值为 $\log4$；第\ref{ch:rl-pomdp}章连续两次相同观测的左状态后验为 $0.6(0.8)^2/[0.6(0.8)^2+0.4(0.2)^2]=24/25$。

\section{证明题与拓展路线}

证明题首先写出对象、假设与所需结论，再确定使用条件期望、压缩性、单调性还是有限求和重排。对于随机算法，必须区分“平均更新的固定点正确”与“随机迭代收敛”；对于近似算法，必须区分“模型内最优”与“真实环境最优”。若结论需要额外覆盖、可积性或平稳性，应把它明确补入命题，而不是用实验曲线代替证明。

理论方向可继续学习马尔可夫链、鞅与随机逼近、统计学习中的集中界和最优控制；算法方向可深入离策略估计、表示学习、策略优化与模型规划；应用方向应加强任务建模、系统辨识和实验设计。基础教材\cite{sutton-barto}适合系统复习，专题原始论文适合核对具体目标、算法版本和理论条件。阅读进阶文献前，能独立复现一个小问题的精确解与采样近似，是检验基础是否稳固的有效标准。



